\documentclass[11pt]{article}
\usepackage[utf8]{inputenc}
\usepackage[T1]{fontenc}
\usepackage{lmodern}
\usepackage[margin=1in]{geometry}
\usepackage{booktabs}
\usepackage{tabularx}
\usepackage{enumitem}
\usepackage{amsmath}
\usepackage[numbers,sort&compress]{natbib}
\usepackage{url}
\usepackage[hidelinks]{hyperref}
\setlist{itemsep=2pt,topsep=4pt}
\newcommand{\file}[1]{\texttt{#1}}

\title{Reconciling Top8000 Ramachandran data with the Richardson Laboratory \texttt{rama8000} reference distributions}
\author{Michael G. Prisant}
\date{15 October 2026}

\begin{document}
\maketitle

\section*{Background}

The MolProbity Ramachandran criteria \cite{Ramachandran1963,Lovell2003,Chen2010} were updated
from the Top8000, a set of about 8000 high-resolution, quality-filtered protein chains, using six
residue categories (General, Gly, Ile/Val, pre-Pro, trans-Pro, cis-Pro)
\cite{Read2011,Richardson2013,Williams2018}. The chains were selected in 2011
\cite{Keedy2011Origins,Top8000Page}; later reference sets were used for rotamers \cite{Lovell2000,Hintze2016}
and for the Top2018 best-parts dataset \cite{Williams2022,Top2018Zenodo}. Cis peptides are discussed in
\cite{Richardson2025}; other non-parametric estimates of Ramachandran distributions are
described in \cite{Ting2010,Shapovalov2013}.

This note reconciles Top8000 Ramachandran data, as we derive them from the laboratory's 2011 chain
files, with the Richardson Laboratory's \file{rama8000} distributions, at the level of category
counts and contour polylines. Our published Top8000 dataset is \cite{Top8000Zenodo}.

\section*{In brief}

\begin{description}[leftmargin=1.5em]
  \item[A.1 Data.] \file{top8000\_chains\_70.tgz}; \file{Top8000\_Rama\_2013.zip}; \file{rama8000}
    NDFT tables; contour Makefile.
  \item[A.2 Programs.] Silk; kin2Dcont; Dangle and Ramalyze (2011 versions) as cross-checks of
    our Python parser.
  \item[B. Implementation.] Revised Ramachandran category precedence, so that pre-Pro Ile/Val go
    into pre-Pro, not Ile/Val; the original Top8000 model files used for atom records; a revised
    rule set that filters on deposited heavy-atom model values only.
\end{description}

\section{Required materials}

These are the original Richardson Laboratory files and programs the comparison needed.

\subsection*{Data files}
\begin{enumerate}
  \item \file{top8000\_chains\_70.tgz} \cite{Top8000Chains}: 7,957 Top8000 chain files from
    2011, with Reduce 3.15 hydrogens \cite{Word1999,Deis2013}. Source: the laboratory website,
    recovered via the Internet Archive. Used as the input for our category counts; the hydrogens
    are ignored.
  \item \file{Top8000\_Rama\_2013.zip} \cite{Top8000Rama2013}: the laboratory's filtered residue lists (one per category)
    plus \file{get\_srcdata.py}, \file{filters\_mysql.txt}, \file{symmetrize.awk}, the Silk
    Makefile and Top8000 Origins \cite{Keedy2011Origins}. Source: laboratory archive. Used as
    the laboratory counts we compared against, and as the source of the filtering rules (the
    SQL in \file{get\_srcdata.py}).
  \item \file{rama8000} NDFT tables (\file{*.ndft}): \file{rlabduke/javadev}, commit
    \file{24e0666} \cite{javadev}. Used as the laboratory grids in the contour comparison, and to
    confirm that the laboratory's filtered residue lists reproduce them exactly.
  \item Contour Makefile: \file{rlabduke/reference\_data}, commit \file{68f581a}
    \cite{referencedata}. Used for the kin2Dcont options and contour levels.
\end{enumerate}

\subsection*{Programs}
\begin{enumerate}[resume]
  \item Silk (two-pass kernel density estimation): \file{rlabduke/javadev}, commit
    \file{24e0666}, rebuilt \cite{javadev}.
  \item kin2Dcont (contour polylines): \file{rlabduke/kincontours}, commit \file{3da22b4},
    rebuilt \cite{kincontours}.
  \item Dangle 1.07 and Ramalyze, mid-2011 versions: \file{rlabduke/javadev}, commit
    \file{eb1c7f48} (2011-06-16), rebuilt \cite{javadev}. Cross-checks only; no reported number
    depends on them.
\end{enumerate}

\paragraph{Not needed.} Reduce and MolProbity. Residue filtering uses deposited heavy-atom model
values only; Reduce hydrogens and flips matter for rotamers, not for Ramachandran filtering.

\paragraph{Not available.} The laboratory's 2011--2013 MySQL database (not preserved);
\file{cont.php}, the contour post-processing script (unpublished).

\paragraph{Our own code.} A Python parser and filters for $\phi$/$\psi$, $\omega$ and the
per-residue values; ContourPy \cite{contourpy} for the marching-squares check.

\section{Comparison}

\subsection{Setup}

\begin{tabularx}{\linewidth}{@{}lX@{}}
  Chain data & (a) The laboratory's 2011 Top8000 chain files (\file{top8000\_chains\_70.tgz}; 7,957
    single chains). (b) Our current RCSB downloads of the same 7,957 chains, as used for our
    published dataset \cite{Top8000Zenodo}. Hydrogens are ignored in both. \\
  Program & Our own Python parser ($\phi$, $\psi$, $\omega$ and per-residue filters). The 2011
    Java Dangle and Ramalyze were run as cross-checks; their $\phi$/$\psi$ agree with ours. \\
  Laboratory reference & The laboratory's filtered residue lists (\file{get\_srcdata.py} output).
    Silk run on these lists reproduces all six \file{rama8000} NDFT tables exactly. \\
\end{tabularx}

\subsection{Chain data: the laboratory's 2011 files and our downloads}

Both start from the chain list \file{Top8000-best\_hom70\_pdb\_chain\_list.csv} \cite{referencedata}
(7,957 chains from 7,737 entries); our copy is byte-identical to the laboratory's. The same list
underlies the published \file{rama8000} contours: the laboratory's filtered residue lists draw on 7,810 of these
chains, 898 of which are not in \file{Top8000-SFbest\_hom70}.

All 7,957 chains (7,737 entries) and their chain IDs agree with the Richardson laboratory's, and we
could download every entry, including three now in the obsolete archive (3m8t, 3n22, 3r4e). For
the mainchain atoms N, CA, C and O, the laboratory's files contain 7,550,002 records and ours
7,550,019. Matching records by chain, residue number, insertion code, atom name and altloc,
7,542,952 are identical in residue name, coordinates, occupancy and $B$. Another 6,150 are
matched but differ in at least one field: residue name 3,264, $B$ 5,605, occupancy 396,
coordinates 4,812, every coordinate change exceeding 1\,\AA. A record can differ in more than one
field. A further 900 records appear only in the laboratory's files and 917 only in ours, from
renumbered, added or removed residues. All 7,967 non-identical records lie in 64 of the 7,737
entries; 8 re-refined or re-sequenced entries (3me5, 3hkx, 3kda, 3oon, 2pvq, 3h87, 2oyp, 1gug)
account for 7,339 of them.

\paragraph{Filtering rules.} All tests use deposited heavy-atom model values only.
\begin{enumerate}
  \item Peptide-linked on both sides (C--N $< 2.0$\,\AA).
  \item Own mainchain (N, CA, C, O): $1 < B < 30$ (the SQL's \file{1 < max\_mc\_b < 30}); for MSE,
    $B < 30$ over all heavy atoms.
  \item Every heavy atom of the residue at occupancy 1.
  \item Non-Pro residues: no alternate location on the preceding residue's C, so that $\phi$ is
    single-valued.
  \item $\omega$ defined; trans $|\omega| \ge 150^\circ$, except cis-Pro $|\omega| < 30^\circ$.
  \item Standard 20 residues + MSE.
  \item General, Ile/Val, pre-Pro only: next residue is standard 20 or MSE.
\end{enumerate}
Rules 3 and 4 are our heavy-atom reading of the SQL's \file{min\_res\_q = 1}; how that column was
computed is not preserved.

\subsection{Category counts}

\begin{center}
\begin{tabular}{lrrrrr}
  \toprule
  & Richardson & \multicolumn{2}{c}{Ours, 2011 files} & \multicolumn{2}{c}{Ours, current files} \\
  \cmidrule(lr){3-4}\cmidrule(lr){5-6}
  Category & laboratory & count & $\Delta$ & count & $\Delta$ \\
  \midrule
  General   & 1,061,639 & 1,061,572 & $-67$ & 1,061,551 & $-88$ \\
  Gly       &   116,589 &   116,579 & $-10$ &   116,573 & $-16$ \\
  Ile/Val   &   195,083 &   195,080 &  $-3$ &   195,083 &   $0$ \\
  pre-Pro   &    61,568 &    61,574 &  $+6$ &    61,573 &  $+5$ \\
  trans-Pro &    66,364 &    66,364 &   $0$ &    66,362 &  $-2$ \\
  cis-Pro   &     3,654 &     3,654 &   $0$ &     3,654 &   $0$ \\
  \midrule
  Total     & 1,504,897 & 1,504,823 & $-74$ & 1,504,796 & $-101$ \\
  \bottomrule
\end{tabular}
\end{center}

On the 2011 files, 90 residues are only in ours and 164 only in the laboratory's lists. On the
current files, 266 are only in ours, 367 only in the laboratory's lists, and 307 have the same
residue number but a different category.

\subsection{Explaining the differences on current files}

Each of the 940 residues that differ on the current files is assigned one cause. If the same rules
on the 2011 files also disagree with the laboratory, it is a rule difference. Otherwise the 2011
file reproduces the laboratory's choice and the cause is the change between the two file versions.

\begin{center}
\begin{tabular}{lrl}
  \toprule
  Cause & Residues & Entries \\
  \midrule
  \textbf{File differences (entries revised since 2011)} & \textbf{686} & \\
  \quad Residue name changed (re-sequenced) & 453 & 3me5, 3hkx, 3h87, 2oyp, +4 \\
  \quad Residue number absent in one file & 160 & 3oon, 3hkx, 2oyp, 3h87, 3me5 \\
  \quad Mainchain $B$ changed & 50 & 2pvq (34), 3hkx, 3me5, +2 \\
  \quad Occupancy, altloc, $\omega$, category, next residue & 23 & 1--5 residues each \\
  \textbf{Rule differences (also on the 2011 files)} & \textbf{254} & 157 entries, at most 14 each \\
  \bottomrule
\end{tabular}
\end{center}

The 254 rule differences are not explained by any heavy-atom rule we have. They include 74 non-Pro
residues that the laboratory kept although the preceding C has an alternate location (rule 4),
out of 10,395 such residues that otherwise pass:

\begin{center}
\begin{tabular}{lrr}
  \toprule
  Preceding C & Laboratory kept & Laboratory dropped \\
  \midrule
  Single altloc record & 37 & 56 \\
  Alternates at the same position & 2 & 43 \\
  Alternates at different positions & 35 & 10,222 \\
  \bottomrule
\end{tabular}
\end{center}

\paragraph{A rule we invented and do not adopt.} Replacing rules 3--4 with ``every atom in the
2011 file at occupancy 1'', which includes the hydrogens added by Reduce 3.15, improves the
agreement on the 2011 files to $+5$ (95 only in ours, 90 only in the laboratory's lists); on the
current files, with Reduce 4.16 hydrogens, it gives $-954$. This rule has no support in the
recollections of the original researchers or in the original filtering artifacts we have, and it
goes against the principle of filtering only on experimentally determined heavy-atom positions.
We report it only as a comparison.

\subsection{Contour polylines}

Contours were drawn with the laboratory's own program, kin2Dcont \cite{kincontours}, using the
options in the laboratory's Makefile \cite{referencedata}, at the Ramalyze levels
\cite{javadev}: inner 98\% for all classes; outer 99.95\% General, 99.80\% cis-Pro, 99.90\%
others.

Laboratory NDFT grid vs our Silk grid, both contoured with kin2Dcont: every contour has the same
number of polylines. The largest deviation is $0.73^\circ$; mean deviations are
$0.00$--$0.07^\circ$.

For comparison, the $1^\circ$ sampling offset between the laboratory's contour grid and the NDFT
grid moves lines by up to $0.88^\circ$, and changing the contouring method moves them by up to
$0.42^\circ$.

\subsubsection*{(i) Density estimation: Silk two-pass adaptive kernel density estimate}

\begin{center}
\small
\begin{tabularx}{\linewidth}{@{}l>{\raggedright\arraybackslash}X>{\raggedright\arraybackslash}X@{}}
  \toprule
  & Original (Richardson laboratory) & Ours \\
  \midrule
  Program & Silk (\file{rlabduke/javadev}) & Silk, javadev \file{24e0666}, rebuilt (OpenJDK 21) \\
  Input & laboratory's filtered residue lists & our $\phi$/$\psi$ under the filtering rules \\
  Gly & mirrored: $(\phi,\psi)$ and $(-\phi,-\psi)$ & same, mirrored explicitly \\
  Kernel, 1st pass & cosine, half-width $5^\circ$ (\file{-cosine=5}) & same \\
  2nd pass & adaptive, \file{-twopass=6.5} (General), 8 (all others) & same \\
  Grid & $2^\circ$ (\file{-gridsize=2}), periodic (\file{-wrap}) & same \\
  Output & \file{-fraction} (fraction of data at lower density) & same \\
  Grid placement & NDFT tables: bins centred on odd degrees (\file{-bounds=-180,180,-180,180});
    contour kinemages: on even degrees (\file{-bounds=-181,179,-181,179})
    & both; contours compared on the odd-degree grid \\
  \bottomrule
\end{tabularx}
\end{center}

Check: Silk with the original options on the laboratory's filtered residue lists reproduces all six
published \file{rama8000} NDFT tables bin for bin.

\subsubsection*{(ii) Contouring}

\paragraph{Original.} kin2Dcont (\file{rlabduke/kincontours}, J.\,M.~Word \cite{kincontours}),
options
\begin{center}\file{-samp -last -nosmooth -group -wrap -180 180 -180 180 -g2}\end{center}
Method: Zyda contour tree
\cite{Zyda1988}, not marching squares.
\begin{itemize}
  \item Each grid value sits on a lattice node; missing samples count as 0.
  \item The wrap edge is duplicated for periodicity.
  \item Each $2^\circ$ cell is split into four triangles around a centre point (the mean of its
    four corners).
  \item Contour crossings are linearly interpolated along triangle edges.
\end{itemize}
Levels (fraction): inner 98\% $=0.02$ for all classes; outer 99.95\% $=0.0005$ General,
99.90\% $=0.001$ others, 99.80\% $=0.002$ cis-Pro. These are the outlier cutoffs in the
laboratory's Ramalyze (\file{Ramachandran.java}; cis-Pro: \file{CISPRO\_ALLOWED}) \cite{javadev};
all are in the Makefile's level list. The laboratory's output then went through an unpublished script (\file{cont.php}).

\paragraph{Ours.}
\begin{itemize}
  \item kin2Dcont, \file{rlabduke/kincontours} \file{3da22b4}, rebuilt with the original options.
    On the odd-degree grid: \file{-wrap -179 181 -179 181}.
  \item Marching squares as a method check: ContourPy \cite{contourpy}, applied to the same grid
    padded periodically by 4 cells.
\end{itemize}

\subsubsection*{(iii) Results}

Laboratory NDFT grid vs our Silk grid, both contoured with kin2Dcont: every contour has the same
number of polylines. The largest deviation is $0.73^\circ$ (Ile/Val outer); mean deviations are
$0.00$--$0.07^\circ$.

For comparison:
\begin{itemize}
  \item the $1^\circ$ offset between the even- and odd-degree grids moves lines by up to
    $0.88^\circ$ (cis-Pro outer, the sparsest class);
  \item kin2Dcont vs marching squares on the same grid differ by up to $0.42^\circ$.
\end{itemize}

Note: the contour comparison used our current RCSB files under the earlier rule set, not the
2011 files and final rules above. So the contour lines are largely insensitive to modest
variation in the residue data set (about 0.1--0.2\%), to the $1^\circ$ grid offset, and to the
contouring method (kin2Dcont vs marching squares). The density estimate itself was not varied:
both sides used the same Silk program and parameters.

\bibliographystyle{unsrtnat}
\bibliography{rama8000_reconciliation}

\end{document}
